\documentclass[11pt,a4paper]{article}
\usepackage[T1]{fontenc}
\usepackage{lmodern}
\usepackage[margin=18mm,top=19mm,bottom=19mm,headheight=14pt,headsep=7mm,footskip=10mm]{geometry}
\usepackage{booktabs,microtype,xcolor,fancyhdr}
\usepackage[hidelinks]{hyperref}
\definecolor{ink}{HTML}{173B4B}
\definecolor{muted}{HTML}{52656C}
\pagestyle{fancy}
\fancyhf{}
\fancyhead[L]{\small\textsc{PhDs but not for long}}
\fancyhead[R]{\small Neural Wave 2026\enspace|\enspace AET}
\fancyfoot[L]{\footnotesize\color{muted}Research report\enspace\textbullet\enspace 18 October 2026}
\fancyfoot[R]{\footnotesize\thepage}
\renewcommand{\headrulewidth}{0.4pt}
\setlength{\parindent}{0pt}
\setlength{\parskip}{5pt}
\setlength{\tabcolsep}{5pt}
\renewcommand{\arraystretch}{1.1}
\newcommand{\reportsection}[2]{\vspace{5pt}{\large\bfseries\color{ink}#1\enspace #2}\par\vspace{2pt}}
\newcommand{\runin}[1]{\textbf{#1}\enspace}
\raggedbottom
\hypersetup{pdftitle={Learning Quarter-Hour Price Corrections},pdfauthor={PhDs but not for long}}

\begin{document}
{\LARGE\bfseries\color{ink}Learning Quarter-Hour Price Corrections\par}
\vspace{4pt}
{\large A joint-zone quantile ensemble\par}
\vspace{6pt}

We reconstruct quarter-hour prices from the supplied hourly curves while
preserving each hourly mean. Our approach learns corrections to linear
interpolation and measures their value through \textbf{MAE / MAE of linear
interpolation}, the ratio used in our figures.

Only a few data findings are needed to motivate the model: movement is largest
near hour boundaries, depends strongly on neighboring-hour ramps, and changes
substantially across months and zones. Negative prices are valid observations.
These findings favor a local correction model with access to historical
shapes and the relationships between zones.

\reportsection{1}{Our ansatz space}

\runin{Learn three corrections around interpolation.}
We begin with linear interpolation between hourly centres and adjust each
hour's four values to preserve its supplied average. Our ansatz uses three
independent values to describe an hour's correction: one shifts the early
half-hour against the late half-hour, and the other two redistribute prices
within those halves. They produce four quarter-hour corrections with zero
average. Adding them to interpolation preserves the hourly mean automatically,
so training can concentrate on the shape that interpolation misses.

We explored different models within this common space. Linear and additive
models capture simple ramp effects, including different responses to upward
and downward price changes. An MLP learns interactions, a temporal CNN captures
neighboring-hour patterns, and a 2D CNN combines those patterns with historical
quarter-hour observations across all four zones.

\runin{The 2D CNN and cross-border relationships.}
For each of the previous seven days, we align quarter-hour prices with hourly
prices repeated at quarter-hour resolution. These form paired rows in a
two-dimensional input. The current day's known hourly curve is included too;
its unobserved quarter-hour row is masked. Small convolutional filters read
the paired representations over short time windows, and a fusion layer
combines information across days.

All four zones enter jointly as channels. This allows the network to learn
a zone's within-hour shape in the context of movements and price differences
across the other zones. Such cross-border relationships are a useful modeling
assumption: Zones 0 and 1, for example, have strongly correlated deviation
shapes, with a correlation of 0.84. Cross-border information enters through
the supplied price curves and their differences; the model can learn when
a local ramp accompanies a broader movement or a separation between zones.
The complete current-day curves are available at inference.

The CNN also receives analog corrections retrieved from strictly earlier
days with similar hourly curves, using up to 90 days of history. These summarize
how interpolation missed under comparable conditions. This gives the model
both a view of recent raw history and examples selected by similarity.

\runin{Combine complementary corrections.}
The final ensemble averages an MLP, a temporal CNN, and the analog 2D CNN,
with two seeds for each family. Each of the six networks has equal weight.
The combination can reduce errors when its members miss in different directions;
its weights remain fixed for all zones and hours. Linear and additive models
were competitive controls, while the ensemble obtained the lowest observed
mean monthly ratio. More elaborate sequence models did not improve that result.

\runin{Learn uncertainty alongside the correction.}
The models predict a central curve and eleven ordered quantiles. Training
combines absolute error on the central curve with pinball loss on the quantiles.
The central prediction preserves hourly means; the surrounding marginal
quantiles describe uncertainty. A separate calibration block sets the
50/80/90/95\% intervals. Fixed physical scaling worked better than additional
normalization: adding RevIN increased the champion's ratio from 0.708 to
0.725 and widened its intervals, so it was excluded from the final model.

{\footnotesize\color{muted}\runin{Implementation and credits.}
Python scientific libraries and PyTorch, with offline CPU inference. Eligible
models use only official data. OpenAI Codex assisted with code, analysis,
and report drafting. Further details and figures:
\url{https://dnd.dsvr.top/index.html#ens}.}

\newpage
\reportsection{2}{Results and validation assumptions}

\runin{The primary metric is MAE / MAE of linear interpolation.}
We compare model and interpolation errors on exactly the same targets.
The interpolation baseline includes the hourly-mean correction described
above. A ratio of 1 means equal error; 0.75 means 25\% less error.
Our figures and comparisons use this ratio so that a difficult, volatile
month can be compared with a quieter one. The headline score is the simple
average of the eight monthly ratios, giving each month equal weight.

\runin{Pooled MAE across expanding-window trials.}
We move forward one month at a time, using larger and larger training windows.
The first trial trains on October--December 2025, validates on January 2026,
and tests on February. The next adds January to training, validates on
February, and tests on March, continuing through September. Training grows,
while validation and testing each retain a one-month horizon in the future.
Our pooled MAE is the average error across all test predictions from these
trials: equivalently, the average of the eight trial MAEs, weighted by the
number of quarter-hour targets evaluated in each month. Testing several months
keeps the assessment from depending on a single unusually difficult period
such as September 2026. Pooled MAE still reflects differences in monthly
difficulty; the average monthly MAE/interpolation-MAE ratio gives each month
equal weight in our main comparison.

The best observed ensemble reaches \textbf{0.707968}, corresponding to
\textbf{29.20\% less error} under this criterion. Its pooled MAE is
4.10 EUR/MWh, compared with 5.81 EUR/MWh for interpolation. Every evaluated
month improves on the baseline.

\begin{center}
\small
\begin{tabular}{@{}lllr@{}}
\toprule
Fitting end & Select/calibrate & Test & MAE ratio \\
\midrule
Dec 2025 & Jan 2026 & Feb 2026 & 0.754 \\
Jan 2026 & Feb 2026 & Mar 2026 & 0.688 \\
Feb 2026 & Mar 2026 & Apr 2026 & 0.751 \\
Mar 2026 & Apr 2026 & May 2026 & 0.704 \\
Apr 2026 & May 2026 & Jun 2026 & 0.661 \\
May 2026 & Jun 2026 & Jul 2026 & 0.692 \\
Jun 2026 & Jul 2026 & Aug 2026 & 0.696 \\
Jul 2026 & Aug 2026 & Sep 2026 & 0.718 \\
\bottomrule
\end{tabular}
\end{center}

Improvements are largest on rising and falling hours, where interpolation
leaves a systematic correction to learn. Flat hours already have small errors
and offer less room for improvement. The calibrated 90\% interval covers
90.6\% of research targets, with mean width 18.8 EUR/MWh. Wider bands generally
coincide with larger errors and greater baseline difficulty; their width is
a relative uncertainty signal, not a probability that a prediction is correct.

\runin{Keep fitting, selection, and calibration chronological.}
Across these trials, fitting grows from 91 to 302 dates, and each model is
fitted from scratch. Days 1--15 of the validation month select settings,
stopping epochs, seed averages, and blend weights; days 16--end calibrate
intervals. The test month is evaluated with parameters and calibration frozen.
All observations from a date remain together, and learned preprocessing uses
fitting data only.

Observed quarter-hour history must strictly precede the reconstructed day.
Earlier test days may supply history for later test days, matching the
challenge's information flow, but they do not trigger model updates.
The daily input still includes all four zones' complete known hourly curves.

\runin{Freeze the delivery recipe separately.}
The research journal records the final three-family recipe refitted on
348 dates, from 1 October 2025 to 15 September 2026. The remaining 15 dates,
16--30 September, are reserved exclusively for interval calibration.
Architecture, weights, learning rates, and epoch budgets are fixed from the
research selection evidence. Inference runs offline without fitting or
recalibration. The scores above belong to the eight research backtests;
they are not a held-out evaluation of this final refit.

\runin{What our validation assumes.}
The procedure assumes older history remains useful and that recent selection
and calibration data are informative about the next month. The dataset
contains only one annual cycle, and dependence between quarters, zones, and
folds limits how much independent evidence it provides. Changing test season
also accompanies increasing training size, so the backtests do not isolate
the benefit of more data.

All months, including September, were inspected during development, and the
champion was chosen after aggregate comparisons. A policy selecting recipes
using only prior data scored 0.7089 in the same campaign. These are
developmental results; small differences between recipes do not establish
superiority on unseen dates. Future interval coverage and the final refit's
accuracy remain assumptions to check on new official observations.

\runin{Conclusion.}
Learning corrections to interpolation lets us focus on quarter-hour shape
while preserving the supplied hourly means. Combining the MLP, temporal CNN,
and cross-border 2D CNN gives a mean monthly MAE/interpolation-MAE ratio of
0.708 in our expanding-window developmental backtests. These results support
combining recent history, similar past days, and relationships between zones.
The frozen ensemble carries these assumptions forward; new official dates
will show how well its accuracy and calibrated uncertainty transfer to future
price regimes.
\end{document}
